\chapter{Markdown as Rhino text}\label{ch:markdowntext}

The command \code{MarkdownText} makes Markdown into Rhino objects: headings, paragraphs, lists, tables and links laid
out as real text, lines and hatches in the model. Double-click any part of it later, and its Markdown opens again for
editing; the same objects are updated in place.

\section{Making Markdown text}\label{sec:mdcreate}

Run \code{MarkdownText}, type or paste Markdown in its dialog, and click where its top left goes, in the active
construction plane. The dialog has:

\begin{xltabular}{\linewidth}{@{}P{0.16\linewidth}L@{}}
\toprule
\ui{Width} & the width text wraps in, in mm; 0: no wrapping \\
\ui{Size}  & the base size in pt: a paragraph's text height; headings are larger \\
\ui{Font}  & the base font \\
\bottomrule
\end{xltabular}

\code{-MarkdownText} takes the Markdown from a file, with the options \emph{Height} and \emph{Width} in model units.
\setting{Markdown parser} chooses what reads the Markdown: Rhino's Markdig, or the plug-in's, as for Panels
(\see{sec:panelmarkdown}).

\section{What the Markdown becomes}

\begin{xltabular}{\linewidth}{@{}P{0.28\linewidth}L@{}}
\toprule
\textsf{\textbf{Markdown}} & \textsf{\textbf{In Rhino}} \\ \midrule
\code{\# Heading}           & larger bold text; the first two levels with a line under them \\
paragraphs                  & wrapped rich text: \textbf{bold}, \emph{italic}, \code{code} in Consolas \\
\code{[text](target)}       & a link: its words in the link color, a line under them (\cref{ch:links}) \\
\code{- item}, \code{1. item} & bullets, numbers \\
\code{- [ ] task}, \code{- [x] done} & task boxes: drawn squares, a check mark in the done ones -- click one to check it \\
\code{> quote}              & a bar along it: a \code{|} on each of its lines \\
code blocks                 & a frame around them \\
tables                      & a grid of cells, the header row bold \\
\code{---}                  & a line \\
\bottomrule
\end{xltabular}

The whole text is one group. How each element looks -- its size, font, color and spacing, the bullets, the task
boxes, the quote's bar, the links -- comes from the file's Markdown style (\see{sec:mdstyle}).

\section{Editing it}\label{sec:mdedit}

\textbf{Double-click} any part of it -- a link too --, or select it and run \code{MarkdownText}, or run the command and
click a part. Its Markdown opens in the dialog:

\begin{itemize}
  \item A part double-clicked or clicked, or selected alone with \keys{Ctrl+Shift+click}, puts the caret at the start
        of its section and shows the other parts' Markdown in gray, editable too. The whole text selected opens it all.
  \item The whole text is selected in Rhino meanwhile.
  \item \ui{OK} updates the \textbf{same objects} in place: their ids, layers and colors stay. Only when the text needs
        more parts than it has are new ones added, and parts it doesn't need any more are removed.
\end{itemize}

\paragraph{Undo and Redo.} In the dialog, \keys{Ctrl+Z} undoes a word at a time (a step ends at a space, a line break,
a tab, a punctuation mark or a pause in the typing), back to the Markdown as it opened; \keys{Ctrl+Y} or
\keys{Ctrl+Shift+Z} redoes. After \ui{OK}, Rhino's \code{Undo} takes the whole edit back -- the objects and their
Markdown as they were -- and \code{Redo} brings it again; a new text is one Undo too.

\shot{14_markdown_text_dialog}{A double-click on the nested list opens the Markdown: that part in black at the caret,
  the other parts in gray. Width, size and font at the bottom.}{fig:mddialog}

\section{The parts belong together}

The parts of a text are tied together by the text's id, stored as user text on each object, not only by the group:
ungroup it, and any part still edits the whole text. A copy (Copy, \keys{Ctrl+C} and \keys{Ctrl+V}) has its own parts:
edited, it parts ways with the original.

\section{Changes made in Rhino are taken in}

Before the dialog opens, what was done to the parts with Rhino's own tools goes into the Markdown:

\begin{description}
  \item[A text's words edited] (Rhino's text editor, or Properties): They are put where they came from, the formatting
        around them kept.
  \item[The whole text moved, turned or scaled:] It stays there.
  \item[All of a block's parts deleted:] Its Markdown is taken out.
\end{description}

A line on top of the dialog, and the command line, say what was taken in.

\section{Task boxes}

\textbf{Click} a task box -- as you click a link (\setting{Markdown links: follow with}) -- and the task is checked or
unchecked: \code{- [ ]} becomes \code{- [x]} in its Markdown, the text is drawn again in place, and Rhino's \code{Undo}
takes it back. A hint by the mouse says what a click does. What was changed on the parts in Rhino is taken in first, as
for an edit. The boxes are drawn -- a square with sharp corners, a check mark in it -- so every box looks the same,
whatever the font; the style can make them characters instead (\see{sec:mdstyle}), and those are clicked alike.

\section{The Markdown style}\label{sec:mdstyle}

How Markdown looks is a \textbf{Markdown style}: a small part of CSS. A rule names elements -- \code{h1} to \code{h6},
\code{p}, \code{li}, \code{task}, \code{blockquote}, \code{pre}, \code{code}, \code{a}, \code{table}, \code{th},
\code{hr}, \code{body} -- and says, for each, its size in em (one base text height: Rhino's text height of a
paragraph), its font, weight, style and color, the space before and after it, the bullet of each list level
(\code{list-style-type}, one character a level), the task boxes (\code{checked: drawn} or a character, their
\code{size}), the quote's bar (\code{bar: "|"}, \code{line} or \code{none}), the code blocks' frame, the tables' grid,
and the links' color and line (\code{text-decoration-thickness}, \code{text-underline-offset}). \code{body}'s font and
color go for everything. No classes, no nesting; what isn't understood is left out and said in \vn{Log}. The default,
\file{styles\textbackslash markdown_rhino.css} next to the plug-in, lists all it can say, with a comment on each line.

\paragraph{One style per Rhino file.} A Rhino file keeps a copy of its style, for all its Markdown texts. A file starts
with \setting{Markdown style for Rhino} (empty: the default) and keeps a copy with its first Markdown text. In
\code{MarkdownText}'s dialog, the line above \ui{Width} says which style the file has:

\begin{xltabular}{\linewidth}{@{}P{0.2\linewidth}L@{}}
\toprule
\ui{Edit style...} & the style in a small editor; what it doesn't understand is shown as you type \\
\ui{Load...}       & a style file's text \\
\ui{Save as...}    & the style as a file of your own, for other Rhino files or the setting \\
\ui{Default}       & miniVivisection's default style \\
\bottomrule
\end{xltabular}

\ui{OK} draws \textbf{all} of the file's Markdown texts with the new style. A file whose style is an untouched copy of a
default -- of any version -- gets the current default when it is opened, its texts drawn again; a style of your own
never changes by itself. V°Text draws with the active Rhino file's style.

\paragraph{The settings.} \setting{Markdown style for Rhino} and \setting{Markdown style for Grasshopper} point to style
files (empty: the defaults); \ui{Edit} beside each opens the file in the same small editor. The defaults are read-only:
\ui{Save as...} makes your own, and the setting then points to it.

\section{Protecting the file: ProtectModel}\label{sec:protectmodel}

The command \code{ProtectModel} is V°Protect (\see{sec:vprotect}) for Rhino files: is the model still the original?
It says the state on the command line, then offers what can be done.

\begin{description}
  \item[In setup,] the files kept -- none set: the Rhino file itself; \emph{AddFiles}, \emph{ClearFiles} -- and where
        their originals go: \emph{OriginalsIn=Archive}, by default \file{(the file's name).protect.zip} next to it
        (\emph{Archive} chooses another), or \emph{OriginalsIn=File}, in the file itself. \emph{Store} stores the
        files' originals now. Everything is kept in the file's document strings.
  \item[Protect] saves the file and keeps it as the original. From then on it is \emph{original} while the file on
        disk is the original and the open model has no changes, \emph{modified} when either isn't so.
  \item[Revert] (asked first) closes the model, its changes dropped, writes the original over the file when that
        changed, and opens it again. \emph{BackToSetup}: not protected any more; the originals stay.
\end{description}

\paragraph{In itself.} A file can't hold its own bytes. So \emph{Protect} saves it once without its original -- those
bytes are the original -- and once more with them. The file on disk is the original when, read, every object, table,
named view and document string but this original is the same: a fingerprint kept with it.

\paragraph{In Markdown text} \code{\{\{protected\}\}} shows \emph{protected} or \emph{not protected}, \code{\{\{modified\}\}}
\emph{modified}, \emph{not modified} or \emph{not tracked}. The texts with these fields are laid out again when the state
changes -- looked at every 1.5\,s while Rhino is idle; that doesn't mark the model changed. A link
\code{[Revert](protect:revert)} reverts.

\section{Links}

Click a link to follow it (\setting{Markdown links: follow with}: or \keys{Ctrl+click}, or off); a
double-click on it edits. Over a link, a hint by the mouse says where it goes and how to follow it (\cref{ch:links}).

\begin{lesson}{the overview}
\lessonfile{00_overview.3dm}: the overview of all examples, and right of it Markdown made into Rhino text. Double-click
any of it to edit; click its links: a named view, a viewport, the V° logo (objects named \emph{Logo}), GH~11's \emph{Ring} component and canvas
view (Grasshopper opens it), lesson 01 (the Script Editor opens it), web pages; click a task box to check it. The file is protected in itself: its
line says \emph{protected} and \emph{not modified}; move the V° logo, and it says \emph{modified}; its link
\emph{Revert to the original} puts the file back.
\end{lesson}
